% !TeX program = lualatex
\documentclass[12pt,aspectratio=169]{beamer}
\usepackage[utf8]{inputenc}
\usepackage[T1]{fontenc}
\usepackage{amsmath,amsfonts,amssymb}
\usepackage{graphicx}
\usepackage{xcolor}
\usepackage{hyperref}
\usepackage{anyfontsize}
\usepackage{lmodern}
\usepackage{longtable}
\usepackage{array}
\usepackage{booktabs}

% ====== EXTREME FONT REDUCTION ======
\renewcommand{\tiny}{\fontsize{3.5}{4.5}\selectfont}
\renewcommand{\scriptsize}{\fontsize{4.5}{5.5}\selectfont}
\renewcommand{\footnotesize}{\fontsize{5.5}{6.5}\selectfont}
\renewcommand{\small}{\fontsize{6.5}{7.5}\selectfont}
\renewcommand{\normalsize}{\fontsize{7.5}{8.5}\selectfont}
\renewcommand{\large}{\fontsize{8.5}{9.5}\selectfont}
\renewcommand{\Large}{\fontsize{9.5}{10.5}\selectfont}
\renewcommand{\LARGE}{\fontsize{10.5}{12}\selectfont}
\renewcommand{\huge}{\fontsize{11.5}{13.5}\selectfont}
\renewcommand{\Huge}{\fontsize{13}{16}\selectfont}

% ====== COLORS ======
\definecolor{redcorrect}{RGB}{200,30,30}
\definecolor{bluecorrect}{RGB}{30,80,200}
\definecolor{greencheck}{RGB}{0,150,50}
\definecolor{lightgray}{RGB}{245,245,245}
\definecolor{darkgray}{RGB}{40,40,40}
\definecolor{softyellow}{RGB}{255,248,200}
\definecolor{deepblue}{RGB}{20,60,150}
\definecolor{darkred}{RGB}{150,20,20}
\definecolor{orangewarning}{RGB}{255,140,0}
\definecolor{correctgreen}{RGB}{0,120,40}
\definecolor{trapcolor}{RGB}{180,80,0}
\definecolor{purpleconcept}{RGB}{120,50,180}
\definecolor{tealhard}{RGB}{0,120,120}

\setbeamercolor{title}{fg=white,bg=redcorrect}
\setbeamercolor{frametitle}{fg=white,bg=redcorrect}
\setbeamercolor{block title}{fg=white,bg=bluecorrect}
\setbeamercolor{block body}{fg=darkgray,bg=lightgray}
\setbeamercolor{normal text}{fg=darkgray}
\usecolortheme{default}

% ====== CUSTOM COMMANDS ======
\newcommand{\correct}[1]{\textcolor{greencheck}{\textbf{\checkmark}} #1}
\newcommand{\keypoint}[1]{\textcolor{deepblue}{\textbf{#1}}}
\newcommand{\hardconcept}[1]{\textcolor{darkred}{\textbf{#1}}}
\newcommand{\examtraps}[1]{\textcolor{trapcolor}{\textbf{⚠ TRAP:}} #1}
\newcommand{\recallbox}[1]{\fcolorbox{deepblue}{blue!5}{\parbox{0.88\textwidth}{\centering \textbf{REMEMBER:} #1}}}
\newcommand{\stepbox}[1]{%
	\begin{center}
		\fcolorbox{darkgray}{white}{%
			\parbox{0.92\textwidth}{\vspace{0.03cm} #1 \vspace{0.03cm}}%
		}
	\end{center}
}
\newcommand{\answerbox}[1]{\fcolorbox{correctgreen}{green!8}{\parbox{0.88\textwidth}{\centering \textcolor{correctgreen}{\textbf{\checkmark\ CORRECT:}} #1}}}
\newcommand{\wronganswer}[1]{\textcolor{redcorrect}{\textbf{✘}} #1}
\newcommand{\trapbox}[1]{\fcolorbox{trapcolor}{orange!8}{\parbox{0.88\textwidth}{\centering \textcolor{trapcolor}{\textbf{⚠ EXAM TRAP:}} #1}}}
\newcommand{\keyrule}[1]{\textcolor{tealhard}{\textbf{🔑}} #1}

\begin{document}
	
	% ================================================================
	% SLIDE 1: TITLE
	% ================================================================
	\begin{frame}
		\centering
		\vspace{0.5cm}
		{\fontsize{18}{22}\selectfont \textbf{CAMS DOMAIN C}}\\[0.15cm]
		{\fontsize{14}{17}\selectfont \textbf{BUILDING AN AML/AFC COMPLIANCE PROGRAM}}\\[0.2cm]
		{\fontsize{9}{11}\selectfont \textit{HARD QUESTIONS \& EXAM TRAPS — 30\% of Exam}}\\[0.1cm]
		\rule{4cm}{0.5pt}\\[0.1cm]
		{\fontsize{7}{9}\selectfont Core Pillars · Three Lines of Defense · Risk Appetite Statement}\\[0.05cm]
		{\fontsize{7}{9}\selectfont EWRA · Risk-Based Approach · Policies/Standards/Procedures}\\[0.05cm]
		{\fontsize{7}{9}\selectfont KYC/CDD/EDD · Transaction Monitoring · KPIs/KRIs}\\[0.05cm]
		{\fontsize{7}{9}\selectfont Board Reporting · Controls Lifecycle · Risk Assessments}
		\vspace{0.3cm}
	\end{frame}
	
	% ================================================================
	% SLIDE 2: CORE PILLARS OF AFC PROGRAM — THE FOUNDATION TRAP
	% ================================================================
	\begin{frame}
		\frametitle{\fontsize{11}{13}\selectfont \textbf{Core Pillars — The Foundation Trap}}
		\begin{block}{\fontsize{7}{9}\selectfont \textbf{The Hard Question — CAMS Exam Favorite}}
			\scriptsize
			What are the core pillars of an effective AML/AFC compliance program?
		\end{block}
		\vspace{0.02cm}
		\answerbox{\large \textbf{1. Risk Assessment. 2. Internal Controls. 3. Independent Testing. 4. Training. 5. Designated Compliance Officer.}}
		\vspace{0.02cm}
		\begin{block}{\fontsize{7}{9}\selectfont \textbf{Natural Explanation — Why You Got This Wrong}}
			\scriptsize
			\begin{itemize}
				\item \wronganswer{Only risk assessment and monitoring} — \textbf{NO.} There are \textbf{5 core pillars}, not just 2.
				\item \textbf{US BSA Requirement:} Every AML program must have:
				\begin{enumerate}
					\item \textbf{Risk Assessment} — Understand your risks
					\item \textbf{Internal Controls} — Policies, procedures, systems
					\item \textbf{Independent Testing} — Audit the program
					\item \textbf{Training} — Educate employees
					\item \textbf{Designated Compliance Officer (MLRO)} — Responsible person
				\end{enumerate}
				\item \textbf{FATF Rec 18:} Requires internal controls, independent audit, and compliance officer.
				\item \textbf{Key principle:} All 5 pillars are \textbf{required}. Missing any one = non-compliant.
			\end{itemize}
		\end{block}
		\vspace{0.02cm}
		\recallbox{\scriptsize \textbf{Key Rule:} 5 pillars = \textbf{Risk Assessment + Internal Controls + Independent Testing + Training + Compliance Officer}.}
	\end{frame}
	
	% ================================================================
	% SLIDE 3: THREE LINES OF DEFENSE — THE ROLE TRAP
	% ================================================================
	\begin{frame}
		\frametitle{\fontsize{11}{13}\selectfont \textbf{Three Lines of Defense — The Role Trap}}
		\begin{block}{\fontsize{7}{9}\selectfont \textbf{The Hard Question — CAMS Exam Favorite}}
			\scriptsize
			Which correctly describes the roles and responsibilities of the three lines of defense?
		\end{block}
		\vspace{0.02cm}
		\answerbox{\large \textbf{1st Line: Business units manage risk. 2nd Line: Compliance oversees. 3rd Line: Internal Audit assures.}}
		\vspace{0.02cm}
		\begin{block}{\fontsize{7}{9}\selectfont \textbf{Natural Explanation — Why You Got This Wrong}}
			\scriptsize
			\begin{itemize}
				\item \wronganswer{Compliance is the first line} — \textbf{NO.} Compliance is \textbf{second line}.
				\item \wronganswer{Internal Audit manages risk} — \textbf{NO.} Audit \textbf{assesses}, doesn't manage.
				\item \textbf{First Line:} Business units (operations, sales, front office). They \textbf{own and manage} risk.
				\item \textbf{Second Line:} Compliance, risk management, MLRO. They \textbf{oversee and challenge} the first line.
				\item \textbf{Third Line:} Internal Audit. They provide \textbf{independent assurance}.
				\item \textbf{Key principle:} 1st = \textbf{do}. 2nd = \textbf{oversee}. 3rd = \textbf{audit}.
			\end{itemize}
		\end{block}
		\vspace{0.02cm}
		\recallbox{\scriptsize \textbf{Key Rule:} 1st = Operations. 2nd = Compliance. 3rd = Audit.}
	\end{frame}
	
	% ================================================================
	% SLIDE 4: MLRO/BSA OFFICER — THE RESPONSIBILITY TRAP
	% ================================================================
	\begin{frame}
		\frametitle{\fontsize{11}{13}\selectfont \textbf{MLRO/BSA Officer — The Responsibility Trap}}
		\begin{block}{\fontsize{7}{9}\selectfont \textbf{The Hard Question — CAMS Exam Favorite}}
			\scriptsize
			What is the primary responsibility of the MLRO/BSA Officer?
		\end{block}
		\vspace{0.02cm}
		\answerbox{\large \textbf{Oversee the AML program and file suspicious transaction reports with the FIU.}}
		\vspace{0.02cm}
		\begin{block}{\fontsize{7}{9}\selectfont \textbf{Natural Explanation — Why You Got This Wrong}}
			\scriptsize
			\begin{itemize}
				\item \wronganswer{Conduct all investigations personally} — \textbf{NO.} The MLRO \textbf{oversees}, doesn't do all the work.
				\item \wronganswer{Make all compliance decisions alone} — \textbf{NO.} The MLRO has authority but works with committees.
				\item \textbf{MLRO (UK/EU) / BSA Officer (US):}
				\begin{itemize}
					\item Designated person responsible for AML compliance
					\item Must have \textbf{authority and resources}
					\item Files SARs/STRs with the FIU
					\item Reports to senior management and board
				\end{itemize}
				\item \textbf{Key principle:} The MLRO is the \textbf{point of contact} for regulators and FIUs.
			\end{itemize}
		\end{block}
		\vspace{0.02cm}
		\recallbox{\scriptsize \textbf{Key Rule:} MLRO = \textbf{oversight} + \textbf{SAR/STR filing}.}
	\end{frame}
	
	% ================================================================
	% SLIDE 5: RISK APPETITE STATEMENT — THE PURPOSE TRAP
	% ================================================================
	\begin{frame}
		\frametitle{\fontsize{11}{13}\selectfont \textbf{Risk Appetite Statement — The Purpose Trap}}
		\begin{block}{\fontsize{7}{9}\selectfont \textbf{The Hard Question — CAMS Exam Favorite}}
			\scriptsize
			What is the purpose of a Risk Appetite Statement (RAS)?
		\end{block}
		\vspace{0.02cm}
		\answerbox{\large \textbf{To define how much risk the institution is willing to accept and guide control implementation.}}
		\vspace{0.02cm}
		\begin{block}{\fontsize{7}{9}\selectfont \textbf{Natural Explanation — Why You Got This Wrong}}
			\scriptsize
			\begin{itemize}
				\item \wronganswer{To eliminate all risk} — \textbf{NO.} Risk cannot be eliminated, only managed.
				\item \wronganswer{To set the budget for compliance} — \textbf{NO.} It's about \textbf{risk tolerance}, not budget.
				\item \textbf{RAS defines:}
				\begin{itemize}
					\item What risks are \textbf{acceptable}
					\item What risks are \textbf{unacceptable}
					\item How to \textbf{measure} risk
					\item When to \textbf{escalate}
				\end{itemize}
				\item \textbf{Key principle:} RAS must be \textbf{communicated} to all employees and \textbf{implemented} through controls.
			\end{itemize}
		\end{block}
		\vspace{0.02cm}
		\recallbox{\scriptsize \textbf{Key Rule:} RAS = \textbf{acceptable risk level} + \textbf{guide for controls}.}
	\end{frame}
	
	% ================================================================
	% SLIDE 6: EWRA — INHERENT vs RESIDUAL RISK TRAP
	% ================================================================
	\begin{frame}
		\frametitle{\fontsize{11}{13}\selectfont \textbf{EWRA — Inherent vs Residual Risk Trap}}
		\begin{block}{\fontsize{7}{9}\selectfont \textbf{The Hard Question — CAMS Exam Favorite}}
			\scriptsize
			What is the difference between inherent risk and residual risk in an EWRA?
		\end{block}
		\vspace{0.02cm}
		\answerbox{\large \textbf{Inherent risk = risk before controls. Residual risk = risk after controls.}}
		\vspace{0.02cm}
		\begin{block}{\fontsize{7}{9}\selectfont \textbf{Natural Explanation — Why You Got This Wrong}}
			\scriptsize
			\begin{itemize}
				\item \wronganswer{Inherent risk is always higher} — \textbf{NO.} It's \textbf{before} controls, so it's usually higher.
				\item \wronganswer{Residual risk is the only risk that matters} — \textbf{NO.} Both matter.
				\item \textbf{Inherent Risk:} Risk without any controls (e.g., all cash transactions are high risk).
				\item \textbf{Control Effectiveness:} How well controls reduce risk (e.g., CDD, monitoring).
				\item \textbf{Residual Risk:} Risk that remains after controls (e.g., cash transactions with controls are medium risk).
				\item \textbf{Key principle:} Controls must reduce residual risk to \textbf{within risk appetite}.
			\end{itemize}
		\end{block}
		\vspace{0.02cm}
		\recallbox{\scriptsize \textbf{Key Rule:} Inherent = \textbf{before} controls. Residual = \textbf{after} controls.}
	\end{frame}
	
	% ================================================================
	% SLIDE 7: RISK-BASED APPROACH — THE DE-RISKING TRAP
	% ================================================================
	\begin{frame}
		\frametitle{\fontsize{11}{13}\selectfont \textbf{Risk-Based Approach — The De-Risking Trap}}
		\begin{block}{\fontsize{7}{9}\selectfont \textbf{The Hard Question — CAMS Exam Favorite}}
			\scriptsize
			What does the risk-based approach require in terms of controls?
		\end{block}
		\vspace{0.02cm}
		\answerbox{\large \textbf{Building controls commensurate with the identified risks.}}
		\vspace{0.02cm}
		\begin{block}{\fontsize{7}{9}\selectfont \textbf{Natural Explanation — Why You Got This Wrong}}
			\scriptsize
			\begin{itemize}
				\item \wronganswer{Applying the same controls to all customers} — \textbf{NO.} That's \textbf{prescriptive}, not risk-based.
				\item \wronganswer{Declining all high-risk customers} — \textbf{NO.} That's \textbf{de-risking}, which FATF discourages.
				\item \textbf{Risk-based approach means:}
				\begin{itemize}
					\item High risk = more controls (EDD, monitoring)
					\item Medium risk = standard controls (CDD)
					\item Low risk = fewer controls (SDD)
				\end{itemize}
				\item \textbf{Key principle:} Controls must be \textbf{proportionate} to risk. NOT \textbf{one-size-fits-all}.
			\end{itemize}
		\end{block}
		\vspace{0.02cm}
		\recallbox{\scriptsize \textbf{Key Rule:} Risk-based = \textbf{proportionate controls}. NOT de-risking.}
	\end{frame}
	
	% ================================================================
	% SLIDE 8: POLICIES vs STANDARDS vs PROCEDURES — THE HIERARCHY TRAP
	% ================================================================
	\begin{frame}
		\frametitle{\fontsize{11}{13}\selectfont \textbf{Policies/Standards/Procedures — The Hierarchy Trap}}
		\begin{block}{\fontsize{7}{9}\selectfont \textbf{The Hard Question — CAMS Exam Favorite}}
			\scriptsize
			What is the difference between policies, standards, and procedures?
		\end{block}
		\vspace{0.02cm}
		\answerbox{\large \textbf{Policies = WHAT and WHY. Standards = MANDATORY REQUIREMENTS. Procedures = HOW.}}
		\vspace{0.02cm}
		\begin{block}{\fontsize{7}{9}\selectfont \textbf{Natural Explanation — Why You Got This Wrong}}
			\scriptsize
			\begin{itemize}
				\item \wronganswer{They are all the same thing} — \textbf{NO.} Different levels of detail.
				\item \textbf{Policies:}
				\begin{itemize}
					\item High-level statements of intent
					\item Example: "We will comply with all AML laws"
					\item \textbf{WHAT} and \textbf{WHY}
				\end{itemize}
				\item \textbf{Standards:}
				\begin{itemize}
					\item Mandatory requirements
					\item Example: "All customers must be risk-rated"
					\item \textbf{MANDATORY REQUIREMENTS}
				\end{itemize}
				\item \textbf{Procedures:}
				\begin{itemize}
					\item Step-by-step instructions
					\item Example: "Step 1: Collect ID. Step 2: Verify..."
					\item \textbf{HOW}
				\end{itemize}
				\item \textbf{Key principle:} Policies set direction; standards set rules; procedures show how to comply.
			\end{itemize}
		\end{block}
		\vspace{0.02cm}
		\recallbox{\scriptsize \textbf{Key Rule:} Policies = WHAT/WHY. Standards = REQUIREMENTS. Procedures = HOW.}
	\end{frame}
	
	% ================================================================
	% SLIDE 9: HORIZON SCANNING — THE EMERGING RISKS TRAP
	% ================================================================
	\begin{frame}
		\frametitle{\fontsize{11}{13}\selectfont \textbf{Horizon Scanning — The Emerging Risks Trap}}
		\begin{block}{\fontsize{7}{9}\selectfont \textbf{The Hard Question — CAMS Exam Favorite}}
			\scriptsize
			What is horizon scanning in the context of AML policies?
		\end{block}
		\vspace{0.02cm}
		\answerbox{\large \textbf{Identifying emerging risks and regulatory changes that could affect the AML program.}}
		\vspace{0.02cm}
		\begin{block}{\fontsize{7}{9}\selectfont \textbf{Natural Explanation — Why You Got This Wrong}}
			\scriptsize
			\begin{itemize}
				\item \wronganswer{Reviewing past incidents only} — \textbf{NO.} Horizon scanning is \textbf{future-oriented}.
				\item \wronganswer{Checking current compliance} — \textbf{NO.} That's monitoring, not horizon scanning.
				\item \textbf{Horizon scanning includes:}
				\begin{itemize}
					\item Regulatory changes (new laws, guidance)
					\item Emerging risks (crypto, AI, new ML methods)
					\item Industry best practices (Wolfsberg, FATF updates)
				\end{itemize}
				\item \textbf{Key principle:} Policies must be \textbf{forward-looking}, not just reactive.
			\end{itemize}
		\end{block}
		\vspace{0.02cm}
		\recallbox{\scriptsize \textbf{Key Rule:} Horizon scanning = \textbf{future risks + regulatory changes}.}
	\end{frame}
	
	% ================================================================
	% SLIDE 10: KYC/CDD/EDD — THE TIMING TRAP
	% ================================================================
	\begin{frame}
		\frametitle{\fontsize{11}{13}\selectfont \textbf{KYC/CDD/EDD — The Timing Trap}}
		\begin{block}{\fontsize{7}{9}\selectfont \textbf{The Hard Question — CAMS Exam Favorite}}
			\scriptsize
			When must CDD be performed?
		\end{block}
		\vspace{0.02cm}
		\answerbox{\large \textbf{Before establishing a relationship, before occasional transactions, when suspicion arises, or when data is doubted.}}
		\vspace{0.02cm}
		\begin{block}{\fontsize{7}{9}\selectfont \textbf{Natural Explanation — Why You Got This Wrong}}
			\scriptsize
			\begin{itemize}
				\item \wronganswer{Only at account opening} — \textbf{NO.} Also for occasional transactions.
				\item \wronganswer{Only when required by regulator} — \textbf{NO.} Must be proactive.
				\item \textbf{CDD Timing:}
				\begin{itemize}
					\item \textbf{Before} establishing relationship (always)
					\item \textbf{Before} occasional transaction (above threshold)
					\item \textbf{When} suspicion arises (overrides SDD)
					\item \textbf{When} doubt about data (re-verify)
				\end{itemize}
				\item \textbf{Key principle:} CDD is \textbf{ongoing}, not one-time.
			\end{itemize}
		\end{block}
		\vspace{0.02cm}
		\recallbox{\scriptsize \textbf{Key Rule:} CDD = \textbf{before relationship + before occasional + when suspicious + when doubted}.}
	\end{frame}
	
	% ================================================================
	% SLIDE 11: EDD — THE SOW/SOF TRAP
	% ================================================================
	\begin{frame}
		\frametitle{\fontsize{11}{13}\selectfont \textbf{EDD — The SOW/SOF Trap}}
		\begin{block}{\fontsize{7}{9}\selectfont \textbf{The Hard Question — CAMS Exam Favorite}}
			\scriptsize
			What is the difference between Source of Wealth and Source of Funds?
		\end{block}
		\vspace{0.02cm}
		\answerbox{\large \textbf{SOW = overall wealth accumulation. SOF = specific funds for the transaction.}}
		\vspace{0.02cm}
		\begin{block}{\fontsize{7}{9}\selectfont \textbf{Natural Explanation — Why You Got This Wrong}}
			\scriptsize
			\begin{itemize}
				\item \wronganswer{They are the same thing} — \textbf{NO.} Different concepts.
				\item \textbf{Source of Wealth (SOW):}
				\begin{itemize}
					\item \textit{How} did you get your money?
					\item Overall wealth picture
					\item Inheritance, business profits, investments
				\end{itemize}
				\item \textbf{Source of Funds (SOF):}
				\begin{itemize}
					\item \textit{Where} did this specific money come from?
					\item Sale of property, salary, business sale
				\end{itemize}
				\item \textbf{Key principle:} BOTH required for EDD (PEPs, high-risk).
			\end{itemize}
		\end{block}
		\vspace{0.02cm}
		\recallbox{\scriptsize \textbf{Key Rule:} SOW = \textbf{overall wealth}. SOF = \textbf{specific funds}. Both required.}
	\end{frame}
	
	% ================================================================
	% SLIDE 12: CUSTOMER RISK ASSESSMENT — THE UPDATING TRAP
	% ================================================================
	\begin{frame}
		\frametitle{\fontsize{11}{13}\selectfont \textbf{Customer Risk Assessment — The Updating Trap}}
		\begin{block}{\fontsize{7}{9}\selectfont \textbf{The Hard Question — CAMS Exam Favorite}}
			\scriptsize
			When should a customer risk assessment be updated?
		\end{block}
		\vspace{0.02cm}
		\answerbox{\large \textbf{When new information emerges, when risk profile changes, or at regular intervals.}}
		\vspace{0.02cm}
		\begin{block}{\fontsize{7}{9}\selectfont \textbf{Natural Explanation — Why You Got This Wrong}}
			\scriptsize
			\begin{itemize}
				\item \wronganswer{Only at onboarding} — \textbf{NO.} Risk is dynamic.
				\item \wronganswer{Only when regulators ask} — \textbf{NO.} Must be proactive.
				\item \textbf{Update triggers:}
				\begin{itemize}
					\item New transaction patterns
					\item PEP status changes
					\item Adverse media hits
					\item Jurisdictional risk changes
					\item Regular reviews (e.g., annually for high-risk)
				\end{itemize}
				\item \textbf{Key principle:} Risk assessments must be \textbf{dynamic} and \textbf{ongoing}.
			\end{itemize}
		\end{block}
		\vspace{0.02cm}
		\recallbox{\scriptsize \textbf{Key Rule:} Customer risk = \textbf{dynamic}. Update when risk changes.}
	\end{frame}
	
	% ================================================================
	% SLIDE 13: TRANSACTION MONITORING — THE SCENARIO TRAP
	% ================================================================
	\begin{frame}
		\frametitle{\fontsize{11}{13}\selectfont \textbf{Transaction Monitoring — The Scenario Trap}}
		\begin{block}{\fontsize{7}{9}\selectfont \textbf{The Hard Question — CAMS Exam Favorite}}
			\scriptsize
			What is the purpose of transaction monitoring scenarios?
		\end{block}
		\vspace{0.02cm}
		\answerbox{\large \textbf{To detect patterns of suspicious activity based on identified money laundering typologies.}}
		\vspace{0.02cm}
		\begin{block}{\fontsize{7}{9}\selectfont \textbf{Natural Explanation — Why You Got This Wrong}}
			\scriptsize
			\begin{itemize}
				\item \wronganswer{To catch all suspicious transactions} — \textbf{NO.} Scenarios are \textbf{targeted}, not comprehensive.
				\item \wronganswer{To replace human judgment} — \textbf{NO.} Scenarios generate alerts for human review.
				\item \textbf{Scenarios:}
				\begin{itemize}
					\item Structured transactions
					\item Round-trip transactions
					\item High-value transfers
					\item Unusual jurisdictions
				\end{itemize}
				\item \textbf{Key principle:} Scenarios must be \textbf{updated} as typologies evolve.
			\end{itemize}
		\end{block}
		\vspace{0.02cm}
		\recallbox{\scriptsize \textbf{Key Rule:} Transaction monitoring = \textbf{targeted scenarios} + \textbf{human review}.}
	\end{frame}
	
	% ================================================================
	% SLIDE 14: KPIs vs KRIs — THE PERFORMANCE TRAP
	% ================================================================
	\begin{frame}
		\frametitle{\fontsize{11}{13}\selectfont \textbf{KPIs vs KRIs — The Performance Trap}}
		\begin{block}{\fontsize{7}{9}\selectfont \textbf{The Hard Question — CAMS Exam Favorite}}
			\scriptsize
			What is the difference between KPIs and KRIs?
		\end{block}
		\vspace{0.02cm}
		\answerbox{\large \textbf{KPIs measure performance. KRIs measure risk.}}
		\vspace{0.02cm}
		\begin{block}{\fontsize{7}{9}\selectfont \textbf{Natural Explanation — Why You Got This Wrong}}
			\scriptsize
			\begin{itemize}
				\item \wronganswer{They are the same thing} — \textbf{NO.} Different purposes.
				\item \textbf{Key Performance Indicators (KPIs):}
				\begin{itemize}
					\item How well is the program working?
					\item Example: \# of SARs filed, % of alerts reviewed on time
				\end{itemize}
				\item \textbf{Key Risk Indicators (KRIs):}
				\begin{itemize}
					\item What is the risk exposure?
					\item Example: \# of high-risk customers, % of customers with missing data
				\end{itemize}
				\item \textbf{Key principle:} Boards need BOTH to make decisions.
			\end{itemize}
		\end{block}
		\vspace{0.02cm}
		\recallbox{\scriptsize \textbf{Key Rule:} KPI = \textbf{performance}. KRI = \textbf{risk}.}
	\end{frame}
	
	% ================================================================
	% SLIDE 15: BOARD REPORTING — THE CONTENT TRAP
	% ================================================================
	\begin{frame}
		\frametitle{\fontsize{11}{13}\selectfont \textbf{Board Reporting — The Content Trap}}
		\begin{block}{\fontsize{7}{9}\selectfont \textbf{The Hard Question — CAMS Exam Favorite}}
			\scriptsize
			What should board reporting on AFC include?
		\end{block}
		\vspace{0.02cm}
		\answerbox{\large \textbf{KPIs, KRIs, regulatory interactions, emerging risks, and control effectiveness.}}
		\vspace{0.02cm}
		\begin{block}{\fontsize{7}{9}\selectfont \textbf{Natural Explanation — Why You Got This Wrong}}
			\scriptsize
			\begin{itemize}
				\item \wronganswer{Only financial metrics} — \textbf{NO.} AFC reporting is about \textbf{risk}, not profit.
				\item \wronganswer{Only regulatory findings} — \textbf{NO.} Board needs \textbf{forward-looking} info too.
				\item \textbf{Board reports should include:}
				\begin{itemize}
					\item \textbf{KPIs/KRIs}: Performance and risk metrics
					\item \textbf{Regulatory interactions}: What did regulators say?
					\item \textbf{Emerging risks}: What's coming?
					\item \textbf{Control effectiveness}: Are controls working?
				\end{itemize}
				\item \textbf{Key principle:} Board needs a \textbf{complete picture} of risk and performance.
			\end{itemize}
		\end{block}
		\vspace{0.02cm}
		\recallbox{\scriptsize \textbf{Key Rule:} Board reports = \textbf{KPIs + KRIs + regulatory + emerging + controls}.}
	\end{frame}
	
	% ================================================================
	% SLIDE 16: FC REPORTING — THE JURISDICTIONAL TRAP
	% ================================================================
	\begin{frame}
		\frametitle{\fontsize{11}{13}\selectfont \textbf{FC Reporting — The Jurisdictional Trap}}
		\begin{block}{\fontsize{7}{9}\selectfont \textbf{The Hard Question — CAMS Exam Favorite}}
			\scriptsize
			What must institutions consider about financial crime reporting across jurisdictions?
		\end{block}
		\vspace{0.02cm}
		\answerbox{\large \textbf{Reporting requirements differ by jurisdiction and must be met separately.}}
		\vspace{0.02cm}
		\begin{block}{\fontsize{7}{9}\selectfont \textbf{Natural Explanation — Why You Got This Wrong}}
			\scriptsize
			\begin{itemize}
				\item \wronganswer{One global report is enough} — \textbf{NO.} Each jurisdiction has its own rules.
				\item \wronganswer{Only home country reporting matters} — \textbf{NO.} Host countries also require reporting.
				\item \textbf{Examples of differences:}
				\begin{itemize}
					\item US: SARs to FinCEN
					\item UK: STRs to NCA
					\item EU: STRs to national FIU
				\end{itemize}
				\item \textbf{Key principle:} Must comply with \textbf{all} jurisdictions where you operate.
			\end{itemize}
		\end{block}
		\vspace{0.02cm}
		\recallbox{\scriptsize \textbf{Key Rule:} Reporting = \textbf{jurisdiction-specific}. Must comply with all.}
	\end{frame}
	
	% ================================================================
	% SLIDE 17: RISK TYPES — THE NON-FINANCIAL RISK TRAP
	% ================================================================
	\begin{frame}
		\frametitle{\fontsize{11}{13}\selectfont \textbf{Risk Types — The Non-Financial Risk Trap}}
		\begin{block}{\fontsize{7}{9}\selectfont \textbf{The Hard Question — CAMS Exam Favorite}}
			\scriptsize
			What types of risk must be managed in an AML program?
		\end{block}
		\vspace{0.02cm}
		\answerbox{\large \textbf{Financial risk, operational risk, non-financial risk, and reputational risk.}}
		\vspace{0.02cm}
		\begin{block}{\fontsize{7}{9}\selectfont \textbf{Natural Explanation — Why You Got This Wrong}}
			\scriptsize
			\begin{itemize}
				\item \wronganswer{Only financial risk} — \textbf{NO.} AML is about \textbf{multiple risk types}.
				\item \textbf{Financial Risk:} Loss from financial crime
				\item \textbf{Operational Risk:} Failure of processes, systems, people
				\item \textbf{Non-Financial Risk:} Compliance, legal, conduct risk
				\item \textbf{Reputational Risk:} Damage to reputation from AML failures
				\item \textbf{Key principle:} All risk types are \textbf{interconnected} in AML.
			\end{itemize}
		\end{block}
		\vspace{0.02cm}
		\recallbox{\scriptsize \textbf{Key Rule:} AML risk = \textbf{financial + operational + non-financial + reputational}.}
	\end{frame}
	
	% ================================================================
	% SLIDE 18: CONTROLS LIFECYCLE — THE ONBOARDING TO EXIT TRAP
	% ================================================================
	\begin{frame}
		\frametitle{\fontsize{11}{13}\selectfont \textbf{Controls Lifecycle — The Onboarding to Exit Trap}}
		\begin{block}{\fontsize{7}{9}\selectfont \textbf{The Hard Question — CAMS Exam Favorite}}
			\scriptsize
			When should AML controls be applied?
		\end{block}
		\vspace{0.02cm}
		\answerbox{\large \textbf{At onboarding, throughout the customer lifecycle, and at exit.}}
		\vspace{0.02cm}
		\begin{block}{\fontsize{7}{9}\selectfont \textbf{Natural Explanation — Why You Got This Wrong}}
			\scriptsize
			\begin{itemize}
				\item \wronganswer{Only at onboarding} — \textbf{NO.} Controls are \textbf{ongoing}.
				\item \wronganswer{Only when suspicious} — \textbf{NO.} Controls are \textbf{continuous}.
				\item \textbf{Controls Lifecycle:}
				\begin{itemize}
					\item \textbf{Onboarding:} CDD, sanctions screening, risk rating
					\item \textbf{Ongoing:} Transaction monitoring, periodic reviews, trigger-based updates
					\item \textbf{Exit:} Account closure, SAR filing, liaison with FIUs
				\end{itemize}
				\item \textbf{Key principle:} Controls are \textbf{continuous}, not one-time.
			\end{itemize}
		\end{block}
		\vspace{0.02cm}
		\recallbox{\scriptsize \textbf{Key Rule:} Controls = \textbf{onboarding + ongoing + exit}.}
	\end{frame}
	
	% ================================================================
	% SLIDE 19: EMPLOYEE DUE DILIGENCE — THE SCREENING TRAP
	% ================================================================
	\begin{frame}
		\frametitle{\fontsize{11}{13}\selectfont \textbf{Employee Due Diligence — The Screening Trap}}
		\begin{block}{\fontsize{7}{9}\selectfont \textbf{The Hard Question — CAMS Exam Favorite}}
			\scriptsize
			What does employee due diligence involve?
		\end{block}
		\vspace{0.02cm}
		\answerbox{\large \textbf{Screening employees against sanctions, PEP lists, and for criminal history.}}
		\vspace{0.02cm}
		\begin{block}{\fontsize{7}{9}\selectfont \textbf{Natural Explanation — Why You Got This Wrong}}
			\scriptsize
			\begin{itemize}
				\item \wronganswer{Only background checks} — \textbf{NO.} More than just criminal checks.
				\item \wronganswer{Only for senior management} — \textbf{NO.} All employees should be screened.
				\item \textbf{Employee Due Diligence includes:}
				\begin{itemize}
					\item Sanctions screening
					\item PEP screening
					\item Criminal background checks
					\item Conflict of interest checks
				\end{itemize}
				\item \textbf{Key principle:} Employees can be a risk vector — screen them.
			\end{itemize}
		\end{block}
		\vspace{0.02cm}
		\recallbox{\scriptsize \textbf{Key Rule:} Employee DD = \textbf{sanctions + PEP + criminal + conflict checks}.}
	\end{frame}
	
	% ================================================================
	% SLIDE 20: RISK ASSESSMENT TYPES — THE COMPREHENSIVE TRAP
	% ================================================================
	\begin{frame}
		\frametitle{\fontsize{11}{13}\selectfont \textbf{Risk Assessment Types — The Comprehensive Trap}}
		\begin{block}{\fontsize{7}{9}\selectfont \textbf{The Hard Question — CAMS Exam Favorite}}
			\scriptsize
			What types of risk assessments are needed in an AML program?
		\end{block}
		\vspace{0.02cm}
		\answerbox{\large \textbf{National Risk Assessment, Entity-wide Risk Assessment, and Customer Risk Assessment.}}
		\vspace{0.02cm}
		\begin{block}{\fontsize{7}{9}\selectfont \textbf{Natural Explanation — Why You Got This Wrong}}
			\scriptsize
			\begin{itemize}
				\item \wronganswer{Only customer risk assessment} — \textbf{NO.} Need \textbf{multiple layers}.
				\item \textbf{Types of Risk Assessments:}
				\begin{itemize}
					\item \textbf{National Risk Assessment (NRA):} Country-level risk
					\item \textbf{Entity-Wide Risk Assessment (EWRA):} Institution-level risk
					\item \textbf{Customer Risk Assessment:} Individual customer risk
					\item \textbf{Product Risk Assessment:} Product-specific risk
					\item \textbf{Jurisdictional Risk Assessment:} Country risk
				\end{itemize}
				\item \textbf{Key principle:} Risk assessment must be \textbf{multi-layered} and \textbf{comprehensive}.
			\end{itemize}
		\end{block}
		\vspace{0.02cm}
		\recallbox{\scriptsize \textbf{Key Rule:} Risk assessments = \textbf{NRA + EWRA + customer + product + jurisdictional}.}
	\end{frame}
	
	% ================================================================
	% SLIDE 21: FINAL EXAM TIPS — DOMAIN C
	% ================================================================
	\begin{frame}
		\frametitle{\fontsize{11}{13}\selectfont \textbf{Domain C — Critical Exam Traps}}
		\scriptsize
		\begin{block}{\fontsize{7}{9}\selectfont \textbf{Most Common Traps — Domain C}}
			\scriptsize
			\begin{enumerate}
				\item \textbf{5 Pillars:} Risk Assessment + Internal Controls + Independent Testing + Training + Compliance Officer.
				\item \textbf{Three Lines:} 1st = Operations. 2nd = Compliance. 3rd = Audit.
				\item \textbf{MLRO:} \textbf{Oversight} + \textbf{SAR/STR filing}. NOT doing all investigations.
				\item \textbf{RAS:} Defines \textbf{acceptable risk} and guides controls. NOT eliminating risk.
				\item \textbf{EWRA:} Inherent = \textbf{before} controls. Residual = \textbf{after} controls.
				\item \textbf{Risk-Based:} \textbf{Proportionate controls}. NOT de-risking.
				\item \textbf{Policies/Standards/Procedures:} Policies = WHAT/WHY. Standards = REQUIREMENTS. Procedures = HOW.
				\item \textbf{Horizon Scanning:} \textbf{Future risks + regulatory changes}. NOT past incidents.
				\item \textbf{CDD Timing:} \textbf{Before relationship + before occasional + when suspicious + when doubted}.
				\item \textbf{SOW vs SOF:} SOW = \textbf{overall wealth}. SOF = \textbf{specific funds}. Both required.
				\item \textbf{Customer Risk:} \textbf{Dynamic}. Update when risk changes.
				\item \textbf{Transaction Monitoring:} \textbf{Targeted scenarios + human review}.
				\item \textbf{KPIs vs KRIs:} KPI = \textbf{performance}. KRI = \textbf{risk}.
				\item \textbf{Board Reporting:} \textbf{KPIs + KRIs + regulatory + emerging + controls}.
				\item \textbf{Reporting:} \textbf{Jurisdiction-specific}. Must comply with all.
				\item \textbf{Risk Types:} \textbf{Financial + operational + non-financial + reputational}.
				\item \textbf{Controls Lifecycle:} \textbf{Onboarding + ongoing + exit}.
				\item \textbf{Employee DD:} \textbf{Sanctions + PEP + criminal + conflict}.
				\item \textbf{Risk Assessments:} \textbf{NRA + EWRA + customer + product + jurisdictional}.
			\end{enumerate}
		\end{block}
		\vspace{0.02cm}
		\recallbox{\scriptsize \textbf{Final Rule:} Domain C = \textbf{program design, governance, and controls}. Understand the \textbf{structure}.}
	\end{frame}
	
	% ================================================================
	% END SLIDE
	% ================================================================
	\begin{frame}
		\centering
		\vspace{1.5cm}
		{\fontsize{18}{22}\selectfont \textbf{Domain C — Complete}}\\[0.2cm]
		{\fontsize{11}{13}\selectfont Building an AML/AFC Compliance Program}\\[0.1cm]
		\rule{4cm}{0.5pt}\\[0.1cm]
		{\fontsize{8}{10}\selectfont Core Pillars · Three Lines of Defense · Risk Appetite Statement}\\[0.05cm]
		{\fontsize{8}{10}\selectfont EWRA · Risk-Based Approach · Policies/Standards/Procedures}\\[0.05cm]
		{\fontsize{8}{10}\selectfont KYC/CDD/EDD · Transaction Monitoring · KPIs/KRIs}\\[0.05cm]
		{\fontsize{8}{10}\selectfont Board Reporting · Controls Lifecycle · Risk Assessments}\\[0.1cm]
		\rule{4cm}{0.5pt}\\[0.1cm]
		{\fontsize{8}{10}\selectfont \textit{All traps explained. All concepts covered. Ready for the exam.}}
		\vspace{0.5cm}
	\end{frame}
	
\end{document}